\section{Maximal almost-disjoint families}

\import{set/almost-disjoint.tex}
\import{set/finite.tex}
\import{order/zorn.tex}
\import{cardinal/infinite.tex}
\import{set/equinumerosity-product.tex}
\import{function.tex}

\begin{definition}\label{maximal_almost_disjoint_family_on}
    $F$ is a maximal almost-disjoint family on $X$ iff
        $F$ is an almost-disjoint family on $X$ and
        for every $G$ such that
            $G$ is an almost-disjoint family on $X$ and
            $F\subseteq G$
        we have $G=F$.
\end{definition}

\begin{theorem}\label{maximal_almost_disjoint_iff_blocking}
    Suppose $F$ is an almost-disjoint family on $X$.
    Then $F$ is a maximal almost-disjoint family on $X$ iff
        $F$ blocks the infinite subsets of $X$.
\end{theorem}
\begin{proof}
    Show that if $F$ is a maximal almost-disjoint family on $X$,
        then $F$ blocks the infinite subsets of $X$.
    \begin{subproof}
        Assume $F$ is a maximal almost-disjoint family on $X$.
        Show that for every $B\subseteq X$ such that $B$ is infinite
            there exists $A\in F$ such that $A\inter B$ is infinite.
        \begin{subproof}
            Fix $B\subseteq X$.
            Assume $B$ is infinite.
            Suppose not.
            For every $A\in F$ we have $A\inter B$ is finite
                by \auto.
            Hence $B$ is almost disjoint from $F$ on $X$
                by \obv{almost_disjoint_from_family_on}.
            Hence $B\notin F$
                by \obv{almost_disjoint_from_family_not_member}.
            $\cons{B}{F}$ is an almost-disjoint family on $X$
                by \auto{almost_disjoint_cons_iff}.
            $F\subseteq\cons{B}{F}$ by \obv{subseteq_cons_self}.
            Hence $\cons{B}{F}=F$
                by \obv{maximal_almost_disjoint_family_on}.
            $B\in\cons{B}{F}$ by \obv{cons_left}.
            Hence $B\in F$.
            Contradiction.
        \end{subproof}
        Follows by \obv{blocks_infinite_subsets_of}.
    \end{subproof}
    Show that if $F$ blocks the infinite subsets of $X$, then
        $F$ is a maximal almost-disjoint family on $X$.
    \begin{subproof}
        Assume $F$ blocks the infinite subsets of $X$.
        Show that for every $G$ such that
            $G$ is an almost-disjoint family on $X$ and
            $F\subseteq G$ we have $G=F$.
        \begin{subproof}
            Fix $G$.
            Assume $G$ is an almost-disjoint family on $X$ and
                $F\subseteq G$.
            Show that every element of $G$ is an element of $F$.
            \begin{subproof}
                Fix $B\in G$.
                \begin{byCase}
                    \caseOf{$B\in F$.}
                        Follows.
                    \caseOf{$B\notin F$.}
                        $B\subseteq X$
                            by \obv{almost_disjoint_members_subset}.
                        $B$ is infinite
                            by \obv{almost_disjoint_members_infinite}.
                        Take $A\in F$ such that $A\inter B$ is infinite
                            by \obv{blocks_infinite_subsets_of}.
                        $A\in G$ by \obv{subseteq}.
                        $A\neq B$.
                        Hence $A\inter B$ is finite by
                            \obv{almost_disjoint_distinct_intersection_finite}.
                        Contradiction.
                \end{byCase}
            \end{subproof}
            Hence $G\subseteq F$ by \obv{subseteq}.
            Follows by \obv{subseteq_antisymmetric}.
        \end{subproof}
        Follows by \obv{maximal_almost_disjoint_family_on}.
    \end{subproof}
    Follows by \auto.
\end{proof}

\begin{corollary}\label{mad_iff_blocking}
    Suppose $F$ is an infinite almost-disjoint family on $X$.
    Then $F$ is an infinite maximal almost-disjoint family on $X$ iff
        $F$ blocks the infinite subsets of $X$.
\end{corollary}
\begin{proof}
    $F$ is an almost-disjoint family on $X$.
    Follows by
        \auto{maximal_almost_disjoint_family_on,
            maximal_almost_disjoint_iff_blocking}.
\end{proof}

\begin{theorem}\label{union_of_chain_of_almost_disjoint_families}
    Suppose $C$ is an inclusion chain.
    Suppose for every $G\in C$ we have
        $G$ is an almost-disjoint family on $X$.
    Then $\unions{C}$ is an almost-disjoint family on $X$.
\end{theorem}
\begin{proof}
    Show that $\unions{C}\subseteq\pow{X}$.
    \begin{subproof}
        It suffices to show that every element of $\unions{C}$ is an element
            of $\pow{X}$ by \obv{subseteq}.
        Fix $A\in\unions{C}$.
        Take $G\in C$ such that $A\in G$ by \obv{unions_iff}.
        $G$ is an almost-disjoint family on $X$.
        Hence $A\subseteq X$
            by \obv{almost_disjoint_members_subset}.
        Follows by \obv{pow_iff,subseteq}.
    \end{subproof}
    Show that for every $A\in\unions{C}$ we have $A$ is infinite.
    \begin{subproof}
        Fix $A\in\unions{C}$.
        Take $G\in C$ such that $A\in G$ by \obv{unions_iff}.
        Follows by \auto{almost_disjoint_members_infinite}.
    \end{subproof}
    Show that for all $A,B\in\unions{C}$ such that $A\neq B$
        we have $A\inter B$ is finite.
    \begin{subproof}
        Fix $A,B\in\unions{C}$.
        Assume $A\neq B$.
        Take $G\in C$ such that $A\in G$ by \obv{unions_iff}.
        Take $H\in C$ such that $B\in H$ by \obv{unions_iff}.
        $G\subseteq H$ or $H\subseteq G$ by \obv{inclusion_chain}.
        \begin{byCase}
            \caseOf{$G\subseteq H$.}
                $A\in H$ by \obv{subseteq}.
                Follows by
                    \auto{almost_disjoint_distinct_intersection_finite}.
            \caseOf{$H\subseteq G$.}
                $B\in G$ by \obv{subseteq}.
                Follows by
                    \auto{almost_disjoint_distinct_intersection_finite}.
        \end{byCase}
    \end{subproof}
    Follows by \obv{almost_disjoint_family_on}.
\end{proof}

\begin{theorem}\label{almost_disjoint_family_has_maximal_extension}
    Suppose $F$ is an almost-disjoint family on $X$.
    Then there exists $M$ such that
        $F\subseteq M$ and
        $M$ is a maximal almost-disjoint family on $X$.
\end{theorem}
\begin{proof}
    Let $P=\{G\in\pow{\pow{X}}\mid
        \text{$F\subseteq G$ and
            $G$ is an almost-disjoint family on $X$}\}$.
    $F\subseteq\pow{X}$ by \obv{almost_disjoint_family_on}.
    Hence $F\in\pow{\pow{X}}$ by \obv{pow_iff,subseteq}.
    $F\subseteq F$ by \obv{subseteq_refl}.
    Hence $F\subseteq F$ and
        $F$ is an almost-disjoint family on $X$.
    Hence $F\in P$.
    Hence $P$ is inhabited.
    Show that for every $C$ such that $C\subseteq P$ and $C$ is inhabited
        and $C$ is an inclusion chain we have $\unions{C}\in P$.
    \begin{subproof}
        Fix $C$.
        Assume $C\subseteq P$ and $C$ is inhabited and
            $C$ is an inclusion chain.
        For every $G\in C$ we have
            $G$ is an almost-disjoint family on $X$
            by \auto{subseteq}.
        Hence $\unions{C}$ is an almost-disjoint family on $X$
            by \obv{union_of_chain_of_almost_disjoint_families}.
        Hence $\unions{C}\subseteq\pow{X}$
            by \obv{almost_disjoint_family_on}.
        Hence $\unions{C}\in\pow{\pow{X}}$
            by \obv{pow_iff,subseteq}.
        Take $G$ such that $G\in C$.
        $G\in P$ by \obv{subseteq}.
        Hence $F\subseteq G$.
        $G\subseteq\unions{C}$ by \auto{unions_iff,subseteq}.
        Hence $F\subseteq\unions{C}$ by \obv{subseteq_transitive}.
    \end{subproof}
    Take $M\in P$ such that for every $G\in P$ such that
        $M\subseteq G$ we have $G=M$ by \obv{zorn_subset_family}.
    $F\subseteq M$ and
        $M$ is an almost-disjoint family on $X$.
    Show that $M$ is a maximal almost-disjoint family on $X$.
    \begin{subproof}
        Show that for every $N$ such that
            $N$ is an almost-disjoint family on $X$ and
            $M\subseteq N$ we have $N=M$.
        \begin{subproof}
            Fix $N$.
            Assume $N$ is an almost-disjoint family on $X$ and
                $M\subseteq N$.
            $N\subseteq\pow{X}$
                by \obv{almost_disjoint_family_on}.
            Hence $N\in\pow{\pow{X}}$ by \obv{pow_iff,subseteq}.
            $F\subseteq N$ by \obv{subseteq_transitive}.
            Hence $N\in P$.
        \end{subproof}
        Follows by \obv{maximal_almost_disjoint_family_on}.
    \end{subproof}
\end{proof}

\begin{theorem}\label{infinite_almost_disjoint_family_has_maximal_extension}
    Suppose $F$ is an infinite almost-disjoint family on $X$.
    Then there exists $M$ such that
        $F\subseteq M$ and
        $M$ is an infinite maximal almost-disjoint family on $X$.
\end{theorem}
\begin{proof}
    $F$ is an almost-disjoint family on $X$.
    Take $M$ such that $F\subseteq M$ and
        $M$ is a maximal almost-disjoint family on $X$
        by \obv{almost_disjoint_family_has_maximal_extension}.
    $F$ is infinite.
    Hence $M$ is infinite by \obv{infinite_superset}.
    $M$ is an almost-disjoint family on $X$
        by \obv{maximal_almost_disjoint_family_on}.
    Hence $M$ is an infinite almost-disjoint family on $X$.
    Hence $M$ is an infinite maximal almost-disjoint family on $X$
        by \obv{maximal_almost_disjoint_family_on}.
\end{proof}

\begin{proposition}\label{natural_pairing_row_subset}
    Let $p\in\bijections{\naturals\times\naturals}{\naturals}$.
    Let $n\in\naturals$.
    Then $\img{p}{\{n\}\times\naturals}\subseteq\naturals$.
\end{proposition}
\begin{proof}
    $p\in\funs{\naturals\times\naturals}{\naturals}$
        by \obv{bijections_to_funs}.
    $\{n\}\subseteq\naturals$ by \obv{singleton_subset_intro}.
    Hence $\{n\}\times\naturals\subseteq
        \naturals\times\naturals$ by \obv{times_subseteq_left}.
    Follows by \obv{funs_image_subseteq}.
\end{proof}

\begin{proposition}\label{natural_pairing_row_infinite}
    Let $p\in\bijections{\naturals\times\naturals}{\naturals}$.
    Let $n\in\naturals$.
    Then $\img{p}{\{n\}\times\naturals}$ is infinite.
\end{proposition}
\begin{proof}
    Let $S=\{n\}\times\naturals$.
    $p\in\funs{\naturals\times\naturals}{\naturals}$ and
        $p$ is injective by \obv{bijections_to_funs,bijections}.
    Hence $p$ is a function and
        $\dom{p}=\naturals\times\naturals$
        by \obv{funs_is_function,funs}.
    For all $x,y\in\naturals\times\naturals$ such that $p(x)=p(y)$
        we have $x=y$ by \auto{injective_function}.
    Hence $p\in\inj{\naturals\times\naturals}{\naturals}$
        by \obv{inj}.
    $\{n\}\subseteq\naturals$ by \obv{singleton_subset_intro}.
    Hence $S\subseteq\naturals\times\naturals$
        by \auto{times_subseteq_left}.
    Let $q=\restrl{p}{S}$.
    $q\in\inj{S}{\naturals}$ by \auto{inj_restrl}.
    $\ran{q}=\img{p}{S}$ by \auto{restrl_ran}.
    Hence $q\in\surj{S}{\img{p}{S}}$ by \auto{inj,fun_to_surj}.
    $q$ is a function and $\dom{q}=S$
        by \auto{inj,funs_is_function,funs}.
    $q$ is injective by \auto{inj,injective_function}.
    Hence $q\in\bijections{S}{\img{p}{S}}$ by \obv{bijections}.
    Hence $S\approx\img{p}{S}$ by \obv{equinum}.
    $S\approx\naturals$ by \auto{times_singleton_left_equinum}.
    Hence $\img{p}{S}\approx\naturals$
        by \obv{equinum_sym,equinum_tran}.
    Suppose not.
    Hence $\img{p}{S}$ is finite.
    Hence $\naturals$ is finite by \obv{finite_equinumerous}.
    Contradiction by \obv{naturals_is_infinite}.
\end{proof}

\begin{proposition}\label{distinct_natural_pairing_rows_disjoint}
    Let $p\in\bijections{\naturals\times\naturals}{\naturals}$.
    Let $n,m\in\naturals$.
    Suppose $n\neq m$.
    Then $\img{p}{\{n\}\times\naturals}\inter
        \img{p}{\{m\}\times\naturals}=\emptyset$.
\end{proposition}
\begin{proof}
    $p\in\funs{\naturals\times\naturals}{\naturals}$ and
        $p$ is injective by \obv{bijections_to_funs,bijections}.
    Hence $p$ is a function and
        $\dom{p}=\naturals\times\naturals$
        by \obv{funs_is_function,funs}.
    Show that every element of
        $\img{p}{\{n\}\times\naturals}\inter
            \img{p}{\{m\}\times\naturals}$
        is an element of $\emptyset$.
    \begin{subproof}
        Fix $z\in
            \img{p}{\{n\}\times\naturals}\inter
                \img{p}{\{m\}\times\naturals}$.
        $z\in\img{p}{\{n\}\times\naturals}$ and
            $z\in\img{p}{\{m\}\times\naturals}$ by \obv{inter}.
        Take $a\in\{n\}\times\naturals$ such that $(a,z)\in p$
            by \obv{img_iff}.
        Take $b\in\{m\}\times\naturals$ such that $(b,z)\in p$
            by \obv{img_iff}.
        Hence $a=b$ by \obv{injective}.
        $\fst{a}=n$ and $\fst{b}=m$
            by \auto{fst_type,singleton_iff}.
        Hence $n=m$.
        Contradiction.
    \end{subproof}
    Follows by \auto{subseteq_emptyset_iff,subseteq}.
\end{proof}

\begin{theorem}\label{infinite_almost_disjoint_family_on_naturals_exists}
    There exists $F$ such that
        $F$ is an infinite almost-disjoint family on $\naturals$.
\end{theorem}
\begin{proof}
    $\naturals$ is an infinite cardinal
        by \obv{naturals_is_infinite,naturals_is_cardinal}.
    Hence $\naturals\times\naturals\approx\naturals$
        by \obv{infinite_cardinal_square}.
    Take $p$ such that
        $p\in\bijections{\naturals\times\naturals}{\naturals}$
        by \obv{equinum}.
    Let $r=\{(n,\img{p}{\{n\}\times\naturals})
        \mid n\in\naturals\}$.
    $r$ is a relation by \auto{relation}.
    $r$ is right-unique by \auto{rightunique,pair_eq_iff}.
    $\dom{r}=\naturals$ by \auto{setext,dom_iff,pair_eq_iff}.
    For every $n\in\naturals$ we have
        $r(n)=\img{p}{\{n\}\times\naturals}$
        by \obv{function_apply_intro}.
    For every $n\in\naturals$ we have $r(n)\in\pow{\naturals}$
        by \auto{natural_pairing_row_subset,pow_iff,subseteq}.
    Hence $r$ is a function to $\pow{\naturals}$.
    Hence $r\in\funs{\naturals}{\pow{\naturals}}$
        by \obv{funs_intro}.
    Show that for all $n,m\in\naturals$ such that $r(n)=r(m)$
        we have $n=m$.
    \begin{subproof}
        Fix $n,m\in\naturals$.
        Assume $r(n)=r(m)$.
        Suppose not.
        Hence $n\neq m$.
        $r(n)=\img{p}{\{n\}\times\naturals}$ and
            $r(m)=\img{p}{\{m\}\times\naturals}$
            by \obv{function_apply_intro}.
        $r(n)$ is infinite by \obv{natural_pairing_row_infinite}.
        Hence $r(n)$ is inhabited by
            \obv{emptyset_is_finite,nonempty_inhabited}.
        Take $z$ such that $z\in r(n)$.
        Hence $z\in r(m)$.
        Hence $z\in r(n)\inter r(m)$ by \obv{inter}.
        $r(n)\inter r(m)=\emptyset$
            by \auto{distinct_natural_pairing_rows_disjoint}.
        Hence $z\in\emptyset$.
        Contradiction by \obv{emptyset}.
    \end{subproof}
    Hence $r\in\inj{\naturals}{\pow{\naturals}}$ by \obv{inj}.
    Let $F=\ran{r}$.
    $F\subseteq\pow{\naturals}$ by \auto{inj,funs_ran}.
    Show that for every $A\in F$ we have $A$ is infinite.
    \begin{subproof}
        Fix $A\in F$.
        Take $n\in\dom{r}$ such that $r(n)=A$ by \obv{fun_ran_iff}.
        Hence $n\in\naturals$ by \obv{funs}.
        Follows by \auto{natural_pairing_row_infinite}.
    \end{subproof}
    Show that for all $A,B\in F$ such that $A\neq B$
        we have $A\inter B$ is finite.
    \begin{subproof}
        Fix $A,B\in F$.
        Assume $A\neq B$.
        Take $n\in\dom{r}$ such that $r(n)=A$ by \obv{fun_ran_iff}.
        Take $m\in\dom{r}$ such that $r(m)=B$ by \obv{fun_ran_iff}.
        Hence $n,m\in\naturals$ by \obv{funs}.
        $n\neq m$.
        Hence $A\inter B=\emptyset$
            by \auto{distinct_natural_pairing_rows_disjoint}.
        Follows by \obv{emptyset_is_finite}.
    \end{subproof}
    Hence $F$ is an almost-disjoint family on $\naturals$
        by \obv{almost_disjoint_family_on}.
    $r\in\surj{\naturals}{F}$ by \auto{inj,fun_to_surj}.
    $r$ is a function and $\dom{r}=\naturals$
        by \obv{inj,funs_is_function,funs}.
    $r$ is injective by \auto{inj,injective_function}.
    Hence $r\in\bijections{\naturals}{F}$ by \obv{bijections}.
    Hence $\naturals\approx F$ by \obv{equinum}.
    Hence $F\approx\naturals$ by \obv{equinum_sym}.
    Show that $F$ is infinite.
    \begin{subproof}
        Suppose not.
        Hence $F$ is finite.
        Hence $\naturals$ is finite by \obv{finite_equinumerous}.
        Contradiction by \obv{naturals_is_infinite}.
    \end{subproof}
    Hence $F$ is an infinite almost-disjoint family on $\naturals$.
\end{proof}

\begin{theorem}\label{infinite_maximal_almost_disjoint_family_on_naturals_exists}
    There exists $M$ such that
        $M$ is an infinite maximal almost-disjoint family on $\naturals$.
\end{theorem}
\begin{proof}
    Take $F$ such that $F$ is an infinite almost-disjoint family on $\naturals$
        by \obv{infinite_almost_disjoint_family_on_naturals_exists}.
    Take $M$ such that $F\subseteq M$ and
        $M$ is an infinite maximal almost-disjoint family on $\naturals$
        by \obv{infinite_almost_disjoint_family_has_maximal_extension}.
\end{proof}
